\section{Sensitivity to the Modeling Assumptions}
\label{sec:robustness}
\subsection{Power Curve and Wake Model}
\label{sec:powercurve}
\label{sec:gaussian}
%% generated by mpce_results.py -- do not edit by hand
\begin{table}[!t]
\centering
\caption{Robustness to the Benchmark Model: Final Layouts Re-Evaluated, Not Re-Optimized}
\label{tab:robust-final}
\scriptsize\setlength{\tabcolsep}{2.4pt}
\begin{tabular}{lcccccc}
\toprule
& \multicolumn{2}{c}{$\Delta$ (\%)} & & & & \\
\cmidrule(lr){2-3}
Model & DS I & DS II & Rank & Best & $\bar\tau$ & Same \\
\midrule
Benchmark & +0.0 & +0.0 & 1.82 & PSO-VNS & -- & -- \\
Cubic curve & -17.7 & -36.9 & 1.82 & PSO-VNS & 0.95 & 94 \\
Cubic, 25 m/s cut-out & -21.5 & -37.0 & 1.80 & PSO-VNS & 0.95 & 93 \\
Gaussian wake & -0.3 & -0.4 & 1.97 & PSO & 0.66 & 54 \\
\bottomrule
\end{tabular}
\par\vspace{2pt}\parbox{\columnwidth}{\scriptsize All feasible final layouts of the 68 cases. $\Delta$: mean relative change of the objective (\%); Rank: average rank of PSO-VNS; Best: best-ranked method; $\bar\tau$: mean Kendall correlation between the benchmark and the alternative ordering within a case; Same: cases (\%) with unchanged best method. Ranks of all methods: supplementary material.}
\end{table}


We re-evaluated every feasible final layout of the main comparison, without re-optimization, with 1) a cubic power curve~\cite{Carrillo2013}, $f_{\rm c}(s)=P_{\text{rated}}(s^3-s_{\text{cut-in}}^3)/(s_{\text{rated}}^3-s_{\text{cut-in}}^3)$ between the cut-in and rated speeds, 2) the same curve with a 25~m/s cut-out, and 3) the Gaussian wake model of Bastankhah and Port\'e-Agel~\cite{Bastankhah2014} with growth rate $K_{\rm G}=0.04$ and the same superposition and Weibull scaling as the benchmark (Section~\ref{S-sec:S-gauss} of the supplementary material).

Table~\ref{tab:robust-final} shows two kinds of sensitivity. First, the absolute energy depends strongly on the model: the cubic curve gives \NCubicOverI\% (Data Set~I) and \NCubicOverII\% (Data Set~II) less expected power than the linear benchmark curve, so the energy values in this paper are benchmark values and not AEP predictions for a specific turbine. Second, the ranking of the algorithms is stable under the cubic curve ($\bar\tau=\NTauCubic$; same best method in \NSameCubic\% of the cases) and less stable under the Gaussian model ($\bar\tau=\NTauGauss$; same best method in \NSameGauss\% of the cases), because layouts optimized for the sharp Jensen wake cone are judged differently by a smooth Gaussian deficit. Under the Gaussian model, PSO-VNS ranks \NPosGaussPSOVNS{} (\NRankGaussPSOVNS), just behind \NBestGauss{} (\NRankGaussPSO), and both remain ahead of all other methods; the small advantage of the VNS phase is thus specific to the model it optimizes.
% CHECK-FINAL [C21]: robustness.Cubic/CubicCutout: tau >= 0.9, PSO-VNS rank position 1
% CHECK-FINAL [C22]: robustness.Gauss: tau < 0.8, same_best < 75, PSOC first, PSOBV second

\subsection{Minimum Spacing and Boundary Handling}
\label{sec:spacing}
\label{sec:boundary}
The $4D$ spacing is inherited from the benchmark literature. A multi-start packing search constructed feasible layouts with up to 13, 26 and 42 turbines at $4D$, 8, 19 and 29 at $5D$, and 7, 13 and 21 at $6D$ for $r=500$, 750 and 1000~m (constructive lower bounds, Table~\ref{S-tab:capacity}), so the 500-m cases with $N\ge9$ (at $5D$) or $N\ge8$ (at $6D$) may be infeasible. The optimized layouts often lie on the $4D$ constraint: of the feasible final layouts of the eight methods, \NSpFiveAll\% also satisfy $5D$ and \NSpSixAll\% satisfy $6D$ (\NSpFiveLarge\% and \NSpSixLarge\% for $N=11$--15), so they cannot be transferred to sites with larger spacing without re-optimization; a re-optimization of six cases at $5D$ and $6D$ with SSA and LX-SSA is reported in Table~\ref{S-tab:spacing-authors} of the supplementary material.
% CHECK-FINAL [C23]: spacing_share.n11_15.pct_5d < 10 ("often lie on the 4D constraint")

The boundary rule also matters: an otherwise identical SSA that projects every turbine radially onto the circle instead of clipping the coordinates attains higher mean benchmark objectives than the original SSA in six representative cases at 3,030 calls, by 26 to 649 units (Mann--Whitney $p<10^{-4}$ in every case). We therefore keep the rule of Section~\ref{sec:constraints} fixed for all methods.

